\begin{table}[htbp]
\centering
\small
\begin{tabular}{p{3.5cm}rrrrrr}
\toprule
Forgatókönyv & Gini(támog.) & CI$_{sz}$ & CI$_{szoc}$ & Kakwani & Theil köz. (\%) & Legr. kvint. (\%) \\
\midrule
0 · Status quo & 0{,}990 & -0{,}725 & -0{,}700 & -1{,}089 & 0 & 0{,}0 \\
1 · Egyenlő fejkvóta & 0{,}407 & -0{,}172 & -0{,}221 & -0{,}536 & 17 & 10{,}5 \\
2 · Epidemiológiai & 0{,}482 & +0{,}191 & +0{,}023 & -0{,}173 & 40 & 22{,}6 \\
3 · Szocioökonómiai & 0{,}416 & +0{,}231 & +0{,}262 & -0{,}134 & 50 & 29{,}6 \\
4 · Kapacitásalapú & 0{,}550 & +0{,}010 & -0{,}154 & -0{,}355 & 23 & 19{,}2 \\
5 · Teljes szükségletalapú & 0{,}498 & +0{,}238 & +0{,}064 & -0{,}126 & 46 & 25{,}3 \\
6 · Költségrés & 0{,}483 & +0{,}193 & +0{,}025 & -0{,}171 & 40 & 22{,}6 \\
7 · Jogszabályi küszöb & 0{,}601 & +0{,}273 & +0{,}324 & -0{,}091 & 23 & 32{,}0 \\
\bottomrule
\end{tabular}
\caption{Méltányossági mérőszámok forgatókönyvenként. CI$_{sz}$: koncentrációs index a szükséglet‑rang szerint (pozitív = progresszív); CI$_{szoc}$: a szocioökonómiai rang szerint; Theil köz.: a vármegyék közötti egyenlőtlenség részaránya.}
\label{tab:equity}
\forras{saját számítás az \texttt{analysis/} pipeline kimeneteiből}
\end{table}
